\section{Energy balls need not be weakly compact in either Banach realization}
\label{sec:r12-energy-ball-weak-compactness-obstruction}

The weak-$*$ compact datum family in (CW14) does not by itself
supply a weakly compact family of candidate potentials. The
following consequence of the proved pole example makes that
obstruction explicit. It is a statement about the two new
Banach norms, not about failure of the complete strong topology.

Fix $n>q\geq1$ and the original radial-mixture background
$S$ of (PO3), on its original unit polydisc $V$. Retain the
original field and absolute value. Put $d=q+1$ and use the
already proved actual potentials
\[
 \psi_j=\max\!\left(\log\max_{1\leq i\leq n}|T_i|,-j\right),
 \qquad
 u=\log\max_{1\leq i\leq n}|T_i|,
                    \qquad j=1,2,\ldots .
 \tag{WB1}
\]
The value of $u$ at the origin is $-\infty$. It represents
a member of the same no-loss first-energy cone, and (PO7)
gives
\[
       e(u)=1,\qquad e(\psi_j)=1-e^{-j},\qquad
                               I(\psi_j,u)=e^{-j}.
 \tag{WB2}
\]
The pole functional is continuous on both $\mathbb D_S$
and $\mathbb L_S$ and satisfies
\[
                   L_\infty(u)=1,\qquad L_\infty(\psi_j)=0.
 \tag{WB3}
\]
All these statements are exactly the earlier original-current
pole construction, including its actual Borel representatives.

\begin{proposition}[No weak cluster and nonreflexivity in every remaining degree]
\label{prop:r12-energy-ball-nonreflexivity}
For the allowed data just specified, the sequence $(\psi_j)$
has no weakly convergent subnet, and in particular no weakly
convergent subsequence, in either $\mathbb D_S$ or
$\mathbb L_S$. Consequently
\[
      \mathcal B_{\mathscr C}
            =\{v\in\mathscr C:r(v)\leq1\}
 \tag{WB4}
\]
is not weakly compact in either Banach space, and neither
$\mathbb D_S$ nor $\mathbb L_S$ is reflexive.

This holds in every remaining degree $q\geq1$ by the
specified example in every $n>q$. It does not say that
these spaces are nonreflexive for every possible background
$S$, including the zero or other degenerate backgrounds.
\end{proposition}

\noindent\textbf{Status: \texttt{DIRECT\_COROLLARY}} of
the present round's exact Banach realizations, source-functional
continuity and pole functional. The proposed universal weak
compactness of bounded potential energy balls is
\texttt{FALSE}. The weak closedness of norm-closed convex
sets and weak compactness of bounded closed balls in a
reflexive Banach space are \texttt{EXTERNAL\_KNOWN}
functional-analysis facts.

\begin{proof}
Let $E$ denote either $\mathbb D_S$ or $\mathbb L_S$.
The proved realization theorems say that the same cone
$\mathscr C$ is norm closed and convex in $E$, and that
\[
                         \|v\|_E=r(v)\qquad(v\in\mathscr C).
 \tag{WB5}
\]
It is therefore weakly closed, by the Hahn--Banach
separation theorem. Every source functional $J(a)$,
$a\in\mathscr C$, is continuous on $E$. In the lattice
case this uses its proved literal extension (OL6),
not a newly assumed extension of an arbitrary dual.
On the cone these source functionals are also continuous
for the strong discrepancy topology, by the frozen
tester estimate and (NI4).

Suppose that a subnet $(\psi_{j_\gamma})$ converges
weakly in $E$ to $v$. The subnet indices are cofinal
in the integer row, so (WB2) still makes it converge
strongly to $u$. Weak closedness gives $v\in\mathscr C$.
For every fixed $a\in\mathscr C$, the weak and strong
continuities just stated give, respectively,
\[
      J(a)(v)=\lim_\gamma J(a)(\psi_{j_\gamma})=J(a)(u).
 \tag{WB6}
\]
Choose $a=v$ and $a=u$. The exact energy-discrepancy
identity then yields
\[
              I(v,u)=(J(v)-J(u))(v-u)=0.
 \tag{WB7}
\]
The strong cone is separated, so $v=u$ as its original
relative Borel class. On the other hand, continuity of
$L_\infty\in E^*$ and (WB3) would give
$L_\infty(v)=0$, whereas $L_\infty(u)=1$. This is a
contradiction. It proves the absence of a weakly
convergent subnet in either space.

The sequence lies in (WB4), by (WB2) and (WB5).
Every net in a compact space has a convergent subnet;
thus (WB4) is not weakly compact. Finally (WB4) is
the intersection of the norm-closed convex cone with
the closed unit ball, so it is weakly closed in that
ball. If $E$ were reflexive, its closed unit ball
would be weakly compact, making (WB4) weakly compact
as well. This contradiction proves nonreflexivity.
\end{proof}

\paragraph{Consequence for the compact-capacity attempt.}
The mechanism ``bounded first energies imply weakly compact
potential candidates; then use weak-$*$ compactness of the
localized source data'' is invalid in these Banach realizations.
The first assertion is disproved by (WB1)--(WB7), even though
the displayed row actually converges strongly. This does not
contradict strong completeness of $\mathscr C$, completeness
of the curvature metric $d_J$, or the weak-$*$ compactness
of the specified datum family in the \emph{dual} space.
It also gives no counterexample to (FG1), to compact ordinary
outer regularity, or to the analytic relation quotient.

\paragraph{Exact dependency snapshots.}
The energy-difference closed-cone theorem is
\path{energy/energy_difference_banach.tex}, SHA256:
\begin{quote}\small
\nolinkurl{12bde9976ef8d8dcb3abbfc23f71f8e1a1ae37a6195de0dcdd13f09c5c0f5f45}.
\end{quote}
The lattice realization and source extension are
\path{nonlinear/order_energy_lattice_completion.tex}, SHA256:
\begin{quote}\small
\nolinkurl{754db6baac08f682e94d5342aa3e757914f3442592bd10cb288f8f2c7ce1bb69}.
\end{quote}
The exact energies and pole dual are
\path{nonlinear/radial_pole_order_dual_obstruction.tex}, SHA256:
\begin{quote}\small
\nolinkurl{f57f4ccac82fc36cbfb4c6663091bdc738d800ebaf022eb9c602176854fa4365}.
\end{quote}
The source strong-continuity and discrepancy identity are
\path{energy/nonlinear_dual_source_geometry.tex}, SHA256:
\begin{quote}\small
\nolinkurl{29a083c0f65d93f777cae9bbe0c02a36385510e7b1a6c9542552322439168743}.
\end{quote}
